% V. Методика за експериментална проверка.
% ВАЖНО: тук няма и не бива да има числа. Разделът описва как ще се мери,
% преди каквото и да е измерване. Числата живеят в Раздел VI.
\section{Методика за експериментална проверка}
\label{sec:method}

\subsection{Цел на експеримента}

Експериментът проверява две твърдения. Първото е, че системата превръща вярно
въпросите от предметната област във формална заявка. Второто е, че перифразата,
върната на потребителя, съответства на действително построената заявка, тоест че
потвърждението от страна на потребителя е информативно, а не формално.

\subsection{Набор от въпроси}

Наборът се съставя предварително и се замразява, преди да е проведено каквото и да е
измерване. Той съдържа въпроси от предметната област, разпределени по вид на
конструкцията, и включва изрично подмножество от нееднозначни въпроси, за които
правилният отговор е повече от едно допустимо тълкуване. Съставът на набора по
категории е даден в \TODO{таблица със състава на набора от въпроси по категории}.

\subsection{Мерки}

Използват се четири мерки. Точност на морфологичния анализ, измерена спрямо ръчно
проверени разчитания. Дял на въпросите, за които правилният извод присъства в гората
от изводи. Дял на въпросите, за които генерираната заявка връща същия резултат като
еталонната заявка при изпълнение върху едни и същи данни. Дял на въпросите, при които
перифразата съответства на построената заявка според ръчна проверка. Пълните
определения на мерките са дадени в \TODO{определения на мерките и на начина на
преброяване}.

\subsection{Съпоставка с базово решение}

Базовото решение е голям езиков модел, на който се подава същият въпрос и същото
описание на схемата, и се иска заявка. То се оценява по същите мерки върху същия
набор. Съпоставката се провежда при фиксирано описание на схемата, за да не се получи
разлика, произтичаща от различен вход. \TODO{конкретният модел, версията му и
настройките, при които е проведена съпоставката}

Отделно от количествения резултат се записва и качественото различие, което е
основната причина за съпоставката: за всеки грешен отговор се отбелязва дали причината
за грешката е установима от изхода на системата. Този показател не се очаква да бъде
симетричен между двете решения.

\subsection{Условия за валидност}

Измерването се смята за валидно само ако наборът от въпроси не е променян след
започване на измерването, ако еталонните заявки са съставени независимо от изхода на
системата, и ако едно и също изпълнение върху едни и същи данни възпроизвежда същите
числа. Всяко нарушение се отбелязва в Раздел~\ref{sec:results} заедно с резултата,
който то засяга.
